\sheet{3}{LTI Systems}{Chapter~3 (Linear Time-Invariant Systems)}

\begin{goals}
\begin{itemize}
  \item Compute convolutions by hand (graphically and in tables) and in C; verify
    the algebraic properties of convolution numerically.
  \item Derive impulse and step responses of the two workhorse smoothers --
    moving average (FIR) and exponential smoother (IIR) -- from their
    difference equations.
  \item Quantify the noise-versus-lag trade-off of a smoother (noise gain,
    mean delay, ramp lag) and relate it to the step response.
  \item Apply LTI filters to the per-tooth engine speed in real time and
    understand why the missing-tooth gap must be handled \emph{before} the
    filter (time invariance in the tooth domain).
\end{itemize}
\end{goals}

\section*{Background}
A discrete-time LTI system is completely described by its impulse response
$h[n]=S\{\delta[n]\}$; the output for any input is the convolution
\[
  y[n]=(x*h)[n]=\sum_{k=-\infty}^{\infty}x[k]\,h[n-k] .
\]
Two smoothers are used throughout the course (written for the per-tooth speed
$r[k]$, but valid for any sequence):
\begin{align*}
  \text{MA (FIR):}\quad & y[n]=\frac1N\sum_{k=0}^{N-1}x[n-k]
     = y[n-1]+\frac{x[n]-x[n-N]}{N},\\
  \text{EMA (IIR):}\quad & y[n]=a\,y[n-1]+(1-a)\,x[n]
     = y[n-1]+\alpha\,(x[n]-y[n-1]),\qquad \alpha=1-a,
\end{align*}
called moving average and exponential smoother (exponential moving average).
Useful characteristics of a smoother with $\sum_n h[n]=1$: the \emph{noise gain}
$\sum_n h^2[n]$ (variance reduction for white noise) and the \emph{mean delay}
$D=\sum_n n\,h[n]$.

%-------------------------------------------------------------------------
\begin{exercise}{Convolution and smoothers by hand}{\Paper}
\begin{tasks}
  \item Compute $y=x*h$ for $x[n]=\{\underline{1},2,0,-1\}$ and
    $h[n]=\{\underline{1},-1,0.5\}$ (underlined: $n=0$) twice: graphically
    (flip and shift $h$) and with a table of products $x[k]h[n-k]$. Check the
    result with $\sum_n y[n]=\sum_n x[n]\cdot\sum_n h[n]$ and explain why this
    identity holds.
  \item Determine impulse response $h[n]$ and step response $s[n]$ of the MA of
    length $N$ and of the EMA (from the difference equation by iteration and
    from $s[n]=\sum_{k\le n}h[k]$). After how many samples does each reach 50\,\%
    and 90\,\% of a step for $N=16$ and $a=15/17$?
  \item Compute noise gain and mean delay $D$ of both smoothers. Show that for a
    ramp input $x[n]=c\,n$ every smoother with $\sum h=1$ has the steady-state
    output $y[n]=c\,(n-D)$, and that $D=\sum_{n\ge0}(1-s[n])$, i.e.\ the ramp lag
    can be read off the step response. Which $a$ gives the EMA the same noise
    gain as the MA of length $N$? Compare the mean delays.
  \item Show that $x[n]=\e^{\jj\Omega n}$ is an eigenfunction of every LTI system,
    $y[n]=H(\Omega)\,x[n]$ with $H(\Omega)=\sum_k h[k]\e^{-\jj\Omega k}$. Evaluate
    $|H|$ for MA$_{16}$, MA$_{30}$ and the EMA with $a=15/17$ at
    $\Omega=2\pi/30$. (On the 60-tooth wheel, sampled once per $6^\circ$, this
    is the 2nd engine order -- the speed ripple caused by the combustions.)
\end{tasks}
\end{exercise}

\begin{solution}
a) Table ($y[n]=\sum_k x[k]h[n-k]$; rows $k$, columns $n$):
\begin{center}
\small
\begin{tabular}{@{}c|rrrrrr@{}}
\toprule
$x[k]\,h[n-k]$ & $n=0$ & 1 & 2 & 3 & 4 & 5 \\
\midrule
$k=0$, $x=1$  & 1 & $-1$ & 0.5 &      &     &   \\
$k=1$, $x=2$  &   & 2    & $-2$ & 1   &     &   \\
$k=2$, $x=0$  &   &      & 0    & 0   & 0   &   \\
$k=3$, $x=-1$ &   &      &      & $-1$& 1   & $-0.5$ \\
\midrule
$y[n]$ & 1 & 1 & $-1.5$ & 0 & 1 & $-0.5$ \\
\bottomrule
\end{tabular}
\end{center}
Graphically: $h[n-k]$ is $h$ flipped and moved to position $n$; $y[n]$ is the
sum of the products where flipped $h$ and $x$ overlap. Length $4+3-1=6$.
Check: $\sum y=1=(1+2+0-1)(1-1+0.5)=2\cdot0.5$. Reason:
$\sum_n\sum_k x[k]h[n-k]=\sum_k x[k]\sum_m h[m]$ (substitute $m=n-k$). The program
of Exercise~3.2 prints \texttt{y = 1 1 -1.5 0 1 -0.5}.

b) MA: $h[n]=1/N$ for $0\le n\le N-1$, else 0;
$s[n]=(n+1)/N$ for $0\le n<N$, $s[n]=1$ for $n\ge N-1$ (linear ramp, finite).
EMA: iterating $y[n]=a y[n-1]+(1-a)\delta[n]$ gives $h[n]=(1-a)a^n u[n]$;
$s[n]=(1-a)\sum_{k=0}^{n}a^k=1-a^{n+1}$ (exponential, never exactly 1).
For $N=16$: MA reaches 50\,\% at $n=7$ ($8/16$), 90\,\% at $n=14$ ($15/16$;
$14/16<0.9$). EMA, $a=15/17$: $a^{n+1}\le0.5\Leftrightarrow n+1\ge\ln0.5/\ln a=5.55$,
i.e.\ $n=5$; 90\,\%: $n+1\ge\ln0.1/\ln a=18.4$, $n=18$. The EMA reacts faster
initially but needs longer to settle.

c) MA: $\sum h^2=N\cdot\frac{1}{N^2}=\frac1N$, $D=\frac1N\sum_{n=0}^{N-1}n=\frac{N-1}{2}$.
EMA: $\sum h^2=(1-a)^2\sum a^{2n}=\frac{(1-a)^2}{1-a^2}=\frac{1-a}{1+a}$,
$D=(1-a)\sum n a^n=\frac{a}{1-a}$.
Ramp: $y[n]=\sum_k h[k]\,c(n-k)=c\,n\sum_k h[k]-c\sum_k k h[k]=c(n-D)$ (for $n$ beyond
the length of $h$; for the EMA asymptotically). With $1-s[n]=\sum_{k>n}h[k]$:
$\sum_{n\ge0}(1-s[n])=\sum_{k}k\,h[k]=D$.
Equal noise: $\frac{1-a}{1+a}=\frac1N\Rightarrow a=\frac{N-1}{N+1}$, $\alpha=\frac{2}{N+1}$
(the usual ``$N$-period EMA''). Then $D_{\text{EMA}}=\frac{(N-1)/(N+1)}{2/(N+1)}=\frac{N-1}{2}=D_{\text{MA}}$:
\emph{for equal white-noise reduction both smoothers have exactly the same ramp lag}.
They differ in the shape of the step response (b) and in how they suppress
periodic disturbances (d). For $N=16$: $a=15/17$, $D=7.5$ samples.

d) $y[n]=\sum_k h[k]\e^{\jj\Omega(n-k)}=\e^{\jj\Omega n}\sum_k h[k]\e^{-\jj\Omega k}$.
MA: geometric sum, $|H|=\left|\frac{\sin(N\Omega/2)}{N\sin(\Omega/2)}\right|$;
EMA: $H=\frac{1-a}{1-a\e^{-\jj\Omega}}$. At $\Omega=2\pi/30$:
MA$_{16}$: $\frac{\sin(16\pi/30)}{16\sin(\pi/30)}=\frac{0.9945}{1.6724}=0.595$;
MA$_{30}$: $\sin(\pi)=0\Rightarrow|H|=0$ -- the MA over one full period removes the
ripple \emph{completely} (the sum over a period of a zero-mean periodic signal
vanishes); EMA: $\frac{0.1176}{|1-0.8824(0.9781-0.2079\jj)|}=\frac{0.1176}{0.2289}=0.514$.
These values are confirmed by the simulation in Exercise~3.3\,a).
\end{solution}

%-------------------------------------------------------------------------
\begin{exercise}{Convolution engine}{\JSLinux}
Program: \codefile{jslinux/sheet03/conv.c}.
\begin{tasks}
  \item Implement \texttt{conv(x, nx, h, nh, y)} (full linear convolution, output
    length $n_x+n_h-1$, no out-of-range accesses) and check it with the example of
    Exercise~3.1\,a).
  \item Verify commutativity, associativity and distributivity with random
    sequences of lengths 200, 50 and 30. Why are the deviations not exactly zero?
    Why might they be exactly zero for ``nice'' test numbers?
  \item Implement the exponential smoother recursively (\texttt{exp\_smooth()}) and
    compare it with the convolution by its impulse response truncated to $L$
    samples ($a=0.9$, $L=10\dots160$). Relate the error to $a^L$ and compare the
    number of multiplications per output sample. Verify the step response
    formula of Exercise~3.1\,b).
  \item Convolve the impulse response of an 8-point MA with itself. Explain the
    shape and the mean delay of the result and compare the step responses of one
    MA and of the cascade.
\end{tasks}
\end{exercise}

\begin{solution}
a) Key part of \codefile{solutions/jslinux/sheet03/conv.c} -- the summation
limits make sure that both $x[k]$ and $h[n-k]$ exist:
\inputcode[firstline=21,lastline=32]{solutions/jslinux/sheet03/conv.c}
Output: \texttt{y = 1 1 -1.5 0 1 -0.5} (as by hand).

b), c) Output:
\begin{lstlisting}[style=shell,numbers=none]
(b) commutativity   max|x*h - h*x|             = 2.66e-15
    associativity   max|(x*h)*g - x*(h*g)|     = 1.07e-14
    distributivity  max|x*(h+g) - (x*h + x*g)| = 5.33e-15

(c) exponential smoother a = 0.90, 1000 samples
      L   max|recursive - conv|   a^L      mult./sample
     10   3.774e-01           3.487e-01   10 (recursion: 2)
     20   1.316e-01           1.216e-01   20 (recursion: 2)
     40   1.600e-02           1.478e-02   40 (recursion: 2)
     80   2.365e-04           2.185e-04   80 (recursion: 2)
    160   5.166e-08           4.773e-08   160 (recursion: 2)
    step response vs. 1 - a^(n+1): max deviation 1.11e-16
\end{lstlisting}
The properties hold exactly in real arithmetic; in floating point the sums are
evaluated in a different order and rounded differently (deviations of a few
$10^{-15}$, i.e.\ some units of the last place relative to values of order 10).
With test numbers that are short binary fractions (e.g.\ multiples of $2^{-15}$)
all products and sums are exact and the deviation is exactly 0 -- a test that
proves nothing about rounding. (c): The truncated impulse response misses the
tail $\sum_{n\ge L}h[n]=a^L$; for an input of magnitude $\approx1.1$ (step + noise)
the error is $\approx1.1\,a^L$, as observed. The recursion needs 2
multiplications per sample for an \emph{infinite} impulse response, the
convolution $L$ -- the reason why IIR filters are cheap. The recursive step
response matches $1-a^{n+1}$ to rounding precision.

d) $h_{\text{MA8}}*h_{\text{MA8}}=\frac{1}{64}\{1,2,\dots,7,8,7,\dots,2,1\}$
(length 15): the overlap of two rectangles of width 8 grows linearly and shrinks
again -- a triangle. Sum 1, mean delay $7=2\cdot3.5$ (mean delays add under
convolution, like the means of independent random variables). The step response
of the cascade is a smooth S-curve (piecewise parabolic) reaching 1 after 15
samples instead of the linear ramp of one MA reaching 1 after 8 samples:
better high-frequency suppression, twice the delay.
\end{solution}

%-------------------------------------------------------------------------
\begin{exercise}{Smoothing the RPM estimate}{\Wokwi}
The per-tooth speed $r[k]=1/T_k$ of Sheet~1 contains the real speed fluctuation
(2nd order, $\pm1\,\%$) and timer quantisation. We smooth it with an MA and an EMA
in the \emph{tooth domain} (one sample per tooth). Three ways to treat the gap
(period $\approx3T$) are compared:
\texttt{GAP\_MODE} 0 -- use $1/T_{\text{gap}}$ like a normal value;
1 -- correct it to $3/T_{\text{gap}}$ (one sample);
2 -- insert $3/T_{\text{gap}}$ three times (one sample per $6^\circ$).
\begin{tasks}
  \item \emph{Offline prediction} (\JSLinux, \codefile{jslinux/sheet03/rpm\_offline.c}).
    The program simulates the crank edges with \SI{0.5}{\micro s} time stamps.
    Implement the gap handling, \texttt{moving\_average()} (running sum) and
    \texttt{exp\_smoother()}. Report: at constant \SI{3000}{rpm} mean, standard
    deviation and min/max of raw, MA$_{16}$, MA$_{30}$, EMA$_{2/17}$ for all
    three gap modes; for a step 3000$\to$\SI{5000}{rpm} the 50\,\%/90\,\% times in
    samples for $N=8,16,32$ and $\alpha=2/(N+1)$; for a ramp of
    \SI{3000}{rpm/s} the lag error in rpm, in samples and in ms.
    Compare with Exercise~3.1.
  \item \emph{Real time} (\Wokwi, \codefile{wokwi/sheet03\_rpmfilter/}). The ISR
    pushes tooth periods into a FIFO; \texttt{loop()} computes and filters the
    speed. Implement the same three steps in the sketch. It prints one line per
    revolution and, whenever the raw speed jumps by more than \SI{150}{rpm} between
    two teeth, a capture block of 128 samples (32 before the jump). Start at
    \SI{3000}{rpm} and \emph{click} the speed slider at \SI{5000}{rpm}. Determine
    the 50\,\%/90\,\% points of both filters from the capture and the ramp lag
    $D=\sum(1-s[n])$. (Dragging the slider produces a staircase of steps, not a
    ramp -- why is the step response sufficient anyway?)
  \item Repeat with \texttt{GAP\_MODE} 0 and 1. Explain the per-revolution values and
    the dips in the capture. Why does the gap violate the time invariance of the
    tooth-domain processing, and why is only mode 2 compatible with the exact
    ripple cancellation of MA$_{30}$?
  \item Design: the injection timing needs the speed during an acceleration of
    \SI{5000}{rpm/s} at \SI{1000}{rpm} with an error below 0.5\,\%, while the
    2nd-order ripple should be suppressed. Is this possible with an MA in the
    tooth domain? What do real ECUs do (keyword: \emph{segment time})?
\end{tasks}
\end{exercise}

\begin{solution}
a) Key parts of \codefile{solutions/jslinux/sheet03/rpm\_offline.c}:
\inputcode[firstline=75,lastline=78]{solutions/jslinux/sheet03/rpm_offline.c}
\inputcode[firstline=92,lastline=112]{solutions/jslinux/sheet03/rpm_offline.c}
Results (constant \SI{3000}{rpm}, after settling):
\begin{lstlisting}[style=shell,numbers=none]
 gap handling: naive gap (2842 samples)
  raw        mean  2966.4  std 260.19  min   996.5  max  3030.3
  MA16       mean  2966.4  std  58.09  min  2859.7  max  3017.8
  MA30       mean  2966.3  std  32.88  min  2932.2  max  3000.0
  EMA 2/17   mean  2966.4  std  57.58  min  2750.7  max  3014.3
 gap handling: gap once (2842 samples)
  raw        mean  3000.3  std  21.53  min  2967.4  max  3030.3
  MA16       mean  3000.3  std  12.19  min  2982.0  max  3017.8
  MA30       mean  3000.4  std   1.06  min  2998.7  max  3002.7
  EMA 2/17   mean  3000.3  std  10.73  min  2984.8  max  3015.5
 gap handling: gap x3 (2940 samples)
  raw        mean  2999.9  std  21.27  min  2967.4  max  3030.3
  MA16       mean  3000.0  std  12.55  min  2982.0  max  3017.8
  MA30       mean  3000.0  std   0.04  min  2999.9  max  3000.1
  EMA 2/17   mean  3000.0  std  10.85  min  2984.8  max  3015.5
\end{lstlisting}
With correct gap handling the raw std ($21.3$\,rpm) is the ripple
($30/\sqrt2$). The attenuation factors $12.55/21.27=0.59$ (MA$_{16}$) and
$10.85/21.27=0.51$ (EMA) agree with $|H|=0.595$ and $0.514$ of Exercise~3.1\,d);
MA$_{30}$ in mode 2 leaves $0.04$\,rpm (timer quantisation only).
Step and ramp (mode 2):
\begin{lstlisting}[style=shell,numbers=none]
  MA8        50% after   4 samples, 90% after   7 samples
  EMA 2/9    50% after   2 samples, 90% after  10 samples
  MA16       50% after   8 samples, 90% after  14 samples
  EMA 2/17   50% after   5 samples, 90% after  18 samples
  MA32       50% after  16 samples, 90% after  28 samples
  EMA 2/33   50% after  11 samples, 90% after  37 samples

  MA16       lag error    6.8 rpm =  7.93 samples =  2.25 ms at 3523 rpm
  EMA 2/17   lag error    6.9 rpm =  7.98 samples =  2.27 ms at 3523 rpm
  MA32       lag error   13.8 rpm = 16.07 samples =  4.56 ms at 3523 rpm
  EMA 2/33   lag error   13.8 rpm = 16.04 samples =  4.55 ms at 3523 rpm
\end{lstlisting}
The step points match 3.1\,b) (MA$_{16}$ is exactly at 50\,\% at $n=7$; the
ripple decides whether 7 or 8 is reported). The ramp lag is
$D+0.5$: $(N-1)/2$ from the filter plus half a sample because $1/T_k$ is the mean
speed over the tooth, i.e.\ the speed half a tooth \emph{before} the edge. MA and
EMA with $\alpha=2/(N+1)$ have the same lag, as predicted in 3.1\,c).

b) Sketch (\codefile{solutions/wokwi/sheet03\_rpmfilter/sketch.ino}), filter
and gap handling:
\inputcode[firstline=67,lastline=82]{solutions/wokwi/sheet03_rpmfilter/sketch.ino}
\inputcode[firstline=129,lastline=138]{solutions/wokwi/sheet03_rpmfilter/sketch.ino}
Expected capture (reference sketch run on the PC with simulated crank edges;
step 3000$\to$5000\,rpm at $k=0$):
\begin{lstlisting}[style=shell,numbers=none]
k,raw,ma,ema,gap        k,raw,ma,ema,gap
-1,3012,3018,3014,0      7,4950,3994,4246,0
0,4988,3143,3247,0       8,4950,4114,4329,0
1,4988,3267,3451,0      14,4988,4846,4671,0
4,4963,3633,3926,0      15,5013,4971,4711,0
5,4950,3753,4046,0      18,5025,4978,4808,0
\end{lstlisting}
The MA rises linearly by $2000/16=125$\,rpm per tooth and reaches the new
value after 16 samples; the EMA jumps by $\alpha\cdot2000=235$\,rpm first and
then approaches exponentially: 50\,\% after 5, 90\,\% after 18 samples
(MA: 8 and 14). The ramp lag follows from the step response, because for an LTI
system the response to a ramp is the running sum of the step response:
$D=\sum_{k\ge0}(1-s[k])$ evaluated on the capture (baseline = mean of the 32
pre-trigger samples, final value = mean of the last 36) gives 7.65 (MA) and
7.56 (EMA) samples, close to 7.5; the extra half sample of the offline ramp
(7.9) is hidden here because the time origin is the first \emph{measured} jump. A staircase is a sum of shifted
steps -- by superposition its response is the sum of shifted step responses;
nothing new is learned from it. Note: the per-revolution line samples the filters
always at the same crank angle (at the gap), so the ripple appears there as a
constant offset (e.g.\ MA$_{16}$ 2982, EMA 2986 at \SI{3000}{rpm}) --
stroboscopic sampling of a crank-synchronous signal.

c) Mode 0: the gap value is $n/3$, i.e.\ a negative pulse of $-2n/3=-2000$\,rpm
once per revolution. MA$_{16}$ shows a rectangular dip of $2000/16=125$\,rpm
lasting 16 teeth (min 2860), the EMA a dip of $\alpha\cdot2000=235$\,rpm (min
2751) that decays over $\approx20$ teeth; the mean is biased by $-1.1\,\%$
(compare 2.1\,c). The per-revolution line, printed right after the gap, shows the
worst value (\SI{3000}{rpm}: MA 2860, EMA 2751; \SI{7000}{rpm}: 6674, 6418).
Mode 1 removes the bias, but the sample covering $18^\circ$ has the same weight as
those covering $6^\circ$; MA$_{30}$ then spans $30$ samples $\neq180^\circ$ for
half of the windows: residual std 1.06\,rpm. Mode 2 makes the sample index
proportional to the crank angle ($6^\circ$ per sample, 60 per revolution); MA$_{30}$
spans exactly $180^\circ$ and cancels the 2nd order completely (std 0.04).
Time invariance: in the tooth domain the operator ``tooth period $\to$ speed''
must be the same for every index $k$. With the gap it is not: at $k\equiv57$
(mod 58) the measurement covers $18^\circ$ instead of $6^\circ$, i.e.\ the system
changes periodically with $k$ (like S1 or S10 of Sheet~2). A shift of the speed
profile by one tooth does not shift the output by one sample. An LTI filter after a
time-variant measurement produces artefacts that no filter design can remove;
the correction must restore a uniform (angle-equidistant) sequence first.

d) At \SI{1000}{rpm} one sample takes \SI{1}{ms}; \SI{5000}{rpm/s} are 5\,rpm per
sample. 0.5\,\% = 5\,rpm allow a lag of only 1 sample, i.e.\ $(N-1)/2+0.5\le1$,
$N\le2$ -- no ripple suppression at all ($|H|\approx1$). Conversely MA$_{30}$
removes the ripple but lags $15.5$ samples $=\SI{15.5}{ms}$, i.e.\ 78\,rpm
($7.8\,\%$). A causal smoother cannot have both. Real ECUs therefore (i) measure
the \emph{segment time} (time for $720^\circ/z = 180^\circ$, i.e.\ exactly the
MA$_{30}$ in the period domain) for torque/misfire functions, synchronous to the
firing, (ii) use the short tooth time for angle extrapolation between teeth,
and (iii) compensate the known lag by prediction ($\hat n=n_{\text{filt}}+\dot n
\cdot D$, an estimate of the acceleration). Note also that the lag in
\emph{samples} is constant while the lag in \emph{time} scales with $1/n$: at
\SI{6000}{rpm} the same MA$_{30}$ lags only \SI{2.6}{ms}.
\end{solution}

\begin{deliverables}
3.1: convolution table, derivations of $h$, $s$, noise gain, $D$, $|H|$. 3.2:
program and output, explanation of the rounding and truncation errors, plot of
the triangular impulse response. 3.3: completed \texttt{rpm\_offline.c} with
result tables; sketch; capture of a step for all three gap modes (plot in
JSLinux), 50/90\,\% points and ramp lag; answer to d).
\end{deliverables}
